\documentclass[12pt,a4paper]{article}

\usepackage[T1]{fontenc}
\usepackage[utf8]{inputenc}
\usepackage[polish]{babel}
\usepackage{lmodern}
\usepackage{enumitem}
\usepackage{geometry}

\geometry{
    a4paper,
    margin=2.5cm
}

\title{\textbf{Topologie arytmetyczne na zbiorze liczb naturalnych}}
\author{}
\date{}

\begin{document}

\maketitle

\begin{center}
\textbf{Propozycja tematyki pracy magisterskiej}
\end{center}

\noindent
\textbf{Kierunek:} Matematyka

\noindent
\textbf{Proponowany promotor:} dr Marta Kwela

\noindent
\textbf{Liczba możliwych prac:} 1--2 (w zależności od wybranego zakresu)

\bigskip

\noindent
\textbf{Opis tematyki.}

Czy własności liczb naturalnych można badać metodami topologicznymi?
Okazuje się, że tak: odpowiednio definiując zbiory bazowe, można wyposażyć
zbiór liczb naturalnych w nietypowe topologie, których bazę stanowią między
innymi nieskończone ciągi arytmetyczne.

Klasycznym przykładem jest topologia Furstenberga, za pomocą której można
uzyskać zaskakujący topologiczny dowód istnienia nieskończenie wielu liczb
pierwszych.

Celem pracy będzie poznanie wybranych topologii na zbiorze liczb naturalnych,
określenie ich podstawowych własności oraz zbadanie, jak arytmetyczna
struktura liczb wpływa na własności topologiczne. W zależności od
zainteresowań i przygotowania studentki lub studenta temat może
koncentrować się na porównaniu topologii Furstenberga, Golomba i Kircha,
na badaniu domknięć i gęstości wybranych zbiorów albo na własnościach
przestrzeni związanych ze spójnością i aksjomatami oddzielania.

Możliwym rozszerzeniem jest włączenie podstaw teorii ideałów -- wprowadzenie
pojęcia zbioru nigdziegęstego i analiza relacji między wybranymi ideałami
topologicznymi. Pozwala to połączyć topologię z teorią mnogości i
kombinatyką oraz porównywać różne sposoby formalnego opisywania
„małych” podzbiorów liczb naturalnych.

Zakres pracy może mieć charakter przede wszystkim przeglądowo-poznawczy,
z możliwością podjęcia prostych samodzielnych problemów i konstrukcji
przykładów.

\bigskip

\noindent
\textbf{Przykładowe zagadnienia do opracowania:}

\begin{itemize}[leftmargin=*]
    \item topologia Furstenberga i topologiczny dowód nieskończoności
    zbioru liczb pierwszych;
    
    \item topologie Golomba i Kircha oraz porównanie ich własności;
    
    \item gęstość, domknięcie i zbiory nigdziegęste w wybranych
    topologiach arytmetycznych;
    
    \item związki między własnościami topologicznymi a podzielnością
    i kongruencjami;
    
    \item ideały zbiorów nigdziegęstych jako narzędzie porównywania
    pojęć „małości” zbiorów.
\end{itemize}

\bigskip

\noindent
\textbf{Literatura podstawowa:}

\begin{enumerate}[leftmargin=*]
    \item H. Furstenberg, \emph{On the infinitude of primes},
    American Mathematical Monthly 62(5) (1955), 353.

    \item S. W. Golomb, \emph{A connected topology for the integers},
    American Mathematical Monthly 66(8) (1959), 663--665.

    \item A. M. Kirch, \emph{A countable, connected, locally connected
    Hausdorff space}, American Mathematical Monthly 76(2) (1969),
    169--171.

    \item L. A. Steen, J. A. Seebach Jr.,
    \emph{Counterexamples in Topology}, wyd.~2,
    Springer-Verlag, 1978.

    \item P. Szyszkowska (Szczuka),
    \emph{Properties of the division topology on the set of positive
    integers}, International Journal of Number Theory 12(3) (2016),
    775--785.
\end{enumerate}

\end{document}
